% SGA 3, Expose VI_B, Loop 1 English reconstruction.
% Scope: Section 3 heading and paragraph 3.0 through Corollary 3.6 in full.
% Authority: Polo--Gille Exp6B-13oct24.pdf, component-local pages 16--18
% (combined-reader pages 357--359; original-margin printed pages 341--343).
% OCR and the Reinhold English Markdown are comparison witnesses only.

\section*{3. 局部有限呈示群的单位连通分支}
\phantomsection\label{sec:sga3-VIB-3}
\addcontentsline{toc}{section}{3. 局部有限呈示群的单位连通分支}

\medskip
\phantomsection\label{par:sga3-VIB-3.0}%
\noindent\textbf{3.0.---}
给定 \(S\)-概形 \(X\) (分别地, \(Y\)) 的子集
\(A\) (分别地, \(B\)), 作为记号滥用, 用
\[
  A\times_S B
\]
表示 \(X\times_S Y\) 中第一个投影属于 \(A\)、
第二个投影属于 \(B\) 的点组成的子集.

给定 \(S\)-群 \(G\) 的子集 \(A\), 若
\[
  c(A)\subset A
  \qquad\text{且}\qquad
  \mu(A\times_S A)\subset A.
\]
则称 \(A\) \emph{在 \(G\) 的群运算下稳定}.

\medskip
\phantomsection\label{def:sga3-VIB-3.1}%
\noindent\textbf{定义 3.1.---}
设 \(G\) 是满足以下条件的取值于群的 \(S\)-函子:
\[
  \tag{+}\label{sga3:vi:definition:3-1:equation:plus}
  \text{对每个 \(s\in S\), 函子
  \(G_s=G\otimes_S\kappa(s)\) 可表.}
\]
用 \(\underline{G_s^0}\) 表示 \(\kappa(s)\)-群 \(G_s\)
的单位元所在的连通分支.%
\footnote{\label{note:sga3-VIB-22}编者注: 这一记号略有滥用,
但当此连通分支是 \(G\) 的开群子概形 \(G^0\) 的底拓扑空间时,
它与第~VI\(_A\)~篇 \SGARefLink{cor:sga3-VIB-2.3}{2.3} 的记号相容;
参见第~VI\(_A\)~篇, \SGARefLink{par:sga3-VIA-2.2bis}{2.2.b}is.}
对每个态射 \(T\to S\), 令
\[
  G^0(T)
    =
  \bigl\{u\in G(T)\bigm|
    \text{对每个 \(s\in S\), }
    u_s(T_s)\subset\underline{G_s^0}\bigr\}.
\]
就定义了 \(G\) 的一个取值于群的 \(S\)-子函子,
称为 \emph{\(G\) 的单位连通分支}, 记为 \(G^0\).
因而定义了函子
\[
  G\longmapsto G^0
\]
从
\[
  \bigl(\widehat{\mathbf{Sch}}_{/S}\bigr)\text{-gr.}
  \quad\text{到}\quad
  \bigl(\widehat{\mathbf{Sch}}_{/S}\bigr)\text{-gr.}
\]

\medskip
\phantomsection\label{rem:sga3-VIB-3.2}%
\noindent\textbf{注 3.2.---}
\textup{(i)} 设 \(G\) 是满足条件 \((+)\) 的取值于群的 \(S\)-函子.
则由第~II~篇,%
\footnote{\label{note:sga3-VIB-23}编者注: 以下为新增内容.}
\[
  \operatorname{Lie}(G^0/S)=\operatorname{Lie}(G/S)
\]
事实上, 对每个 \(S\)-概形 \(T\),
用 \(I_T\) 表示 \(T\) 上的对偶数 \(T\)-概形, 并用
\[
  \tau:T\longrightarrow I_T
\]
表示零截面 (参见 \SGARefLink{sga3:ii:definition:2-1}{第~II~篇, 2.1}).
按照定义,
\[
  \operatorname{Lie}(G/S)(T)
    =
  \{u\in G(I_T)\mid u\circ\tau=e\},
\]
其中 \(e:T\to G\) 表示 \(T\to S\) 与单位截面 \(S\to G\) 的复合.
同样地,
\[
  \begin{aligned}
  \operatorname{Lie}(G^0/S)(T)
    &=
  \{u\in G^0(I_T)\mid u\circ\tau=e\} \\
    &=
  \bigl\{u\in G(I_T)\bigm|
    u\circ\tau=e
    \ \text{且}\
    u_s\bigl((I_T)_s\bigr)\subset\underline{G_s^0}
    \text{ 对每个 \(s\in S\)}\bigr\}.
  \end{aligned}
\]
对每个 \(s\in S\), 概形 \(T_s\) 和
\[
  (I_T)_s=I_{T_s}
\]
具有相同的底集. 因而若 \(u\in\operatorname{Lie}(G/S)(T)\), 则
\[
  u_s(I_{T_s})=e_s,
\]
其中 \(e_s\) 表示 \(G_s\) 的单位点,
从而 \(u\in\operatorname{Lie}(G^0/S)(T)\).
故包含关系
\[
  \operatorname{Lie}(G^0/S)(T)
    \subset
  \operatorname{Lie}(G/S)(T)
\]
对每个 \(T\) 都是等式, 因此
\[
  \operatorname{Lie}(G^0/S)=\operatorname{Lie}(G/S).
\]

\textup{(ii)} 设 \(G\) 和 \(G'\) 是满足 \((+)\) 的两个取值于群的
\(S\)-函子. 则:
\begin{enumerate}
\item[\textup{a)}] 若 \(G\subset G'\), 则 \(G^0\subset G'^0\);
\item[\textup{b)}] 若 \(G\subset G'\) 且 \(G'^0\subset G\), 则 \(G^0=G'^0\);
\item[\textup{c)}] 若对每个 \(s\in S\), \(G_s\)
  在 \(\kappa(s)\) 上局部有限型,
  则由第~VI\(_A\)~篇, \SGARefLink{cor:sga3-VIB-2.3}{2.3},
  \(G^0\) 满足条件 \((+)\), 并且
  \[
    (G^0)^0=G^0.
  \]
\end{enumerate}

\medskip
\phantomsection\label{prop:sga3-VIB-3.3}%
\noindent\textbf{命题 3.3.---}
\emph{设 \(G\) 是满足条件 \((+)\) 的取值于群的 \(S\)-函子,
\(S'\) 是 \(S\)-概形. 则}
\[
  (G\times_S S')^0=G^0\times_S S'.
\]
\emph{换言之, 函子 \(G\mapsto G^0\) 与基变换可交换;
即下图交换:}
\begin{center}
  \phantomsection\label{diag:sga3-VIB-3.3-base-change}%
  \begin{tikzcd}[row sep=large,column sep=huge]
    \bigl(\widehat{\mathrm{Sch}}_{/S}\bigr)\text{-}\mathrm{gr.}
      \arrow[r,"{G\mapsto G^0}"]
      \arrow[d] &
    \bigl(\widehat{\mathrm{Sch}}_{/S}\bigr)\text{-}\mathrm{gr.}
      \arrow[d] \\
    \bigl(\widehat{\mathrm{Sch}}_{/S'}\bigr)\text{-}\mathrm{gr.}
      \arrow[r] &
    \bigl(\widehat{\mathrm{Sch}}_{/S'}\bigr)\text{-}\mathrm{gr.}
  \end{tikzcd}
\end{center}

只须检验, 对每个 \(s'\in S'\), 用 \(s\) 表示它在 \(S\) 中的像, 有
\[
  \begin{aligned}
  \bigl((G\times_S S')\otimes_{S'}\kappa(s')\bigr)^0
    &=
  \bigl(G_s\otimes_{\kappa(s)}\kappa(s')\bigr)^0 \\
    &=
  G_s^0\otimes_{\kappa(s)}\kappa(s').
  \end{aligned}
\]
这由第~VI\(_A\)~篇, \SGARefLink{rem:sga3-VIB-2.1.2}{2.1.2} 得到.
为便于后文 \SGARefLink{prop:sga3-VIB-4.2}{4.2} 使用,
注意这里并未使用 \(G_s\) 的群结构,
而只使用了 \(G_s^0\) 有有理点 \(\epsilon(s)\),
因而几何连通这一事实 (另见 EGA~IV\(_2\), 4.5.14).

\medskip
\phantomsection\label{case:sga3-VIB-3.4}%
\noindent\textbf{特殊情形 3.4.---}
设 \(G\) 为 \(S\)-群概形, 用 \(\underline{G^0}\)
表示 \(G\) 中当 \(s\) 遍历 \(S\) 时各 \(\underline{G_s^0}\)
的并组成的子集.
则 \(\underline{G^0}\) 是 \(G\) 中在 \(G\) 的群运算下稳定的子集
(参见 \hyperref[par:sga3-VIB-3.0]{3.0}), 且对每个态射 \(S'\to S\),
\[
  G^0(S')
    =
  \bigl\{u\in G(S')\bigm|
    u(S')\subset\underline{G^0}\bigr\}.
\]
当 \(\underline{G^0}\) 是 \(G\) 的开子集时,
显然 \(G^0\) 由 \(G\) 在开子集 \(\underline{G^0}\) 上诱导的子概形表示.

\medskip
\phantomsection\label{prop:sga3-VIB-3.5}%
\noindent\textbf{命题 3.5.---}
\emph{设 \(S\) 为拟紧且拟分离的概形,
\(G\) 为纤维局部有限型的 \(S\)-群.
则存在 \(G\) 的拟紧开子集 \(U\), 包含 \(\underline{G^0}\).}

由于单位截面 \(\epsilon\) 为浸入,
\(\epsilon(S)\) 是 \(G\) 的拟紧子空间.
因此存在 \(G\) 的拟紧开子集 \(V\), 包含 \(\epsilon(S)\).
由于 \(S\) 拟分离且 \(V\) 拟紧,
\(V\) 在 \(S\) 上拟紧 (EGA~IV\(_1\), 1.2.4).
因而 \(V\times_SV\) 在 \(S\) 上拟紧, 故拟紧.
由此可知
\[
  V\cdot V=\mu(V\times_SV)
\]
拟紧. 令
\[
  V_s=V\cap G_s,
  \qquad
  V_s^0=V\cap\underline{G_s^0}.
\]
则 \(V_s^0\) 在 \(\underline{G_s^0}\) 中开,
且由于 \(G_s^0\) 不可约
(第~VI\(_A\)~篇, \SGARefLink{cor:sga3-VIB-2.4}{2.4}),
它在 \(\underline{G_s^0}\) 中稠密. 因此
\[
  V_s^0\cdot V_s^0=\underline{G_s^0}
  \qquad
  \text{(第~VI\(_A\)~篇, 0.5),}
\]
这表明
\[
  V_s\cdot V_s\supset\underline{G_s^0},
  \qquad\text{从而}\qquad
  V\cdot V\supset\underline{G^0}.
\]
最后, 由于 \(V\cdot V\) 拟紧,
存在 \(G\) 的拟紧开子集 \(U\), 包含 \(V\cdot V\),
因而也包含 \(\underline{G^0}\).

\medskip
\phantomsection\label{cor:sga3-VIB-3.6}%
\noindent\textbf{推论 3.6.---}
\emph{设 \(G\) 为纤维局部有限型且连通的 \(S\)-群.
则 \(G\) 在 \(S\) 上拟紧.}

此断言在 \(S\) 上为局部问题 (EGA~I, 6.6.1),
故可化归到 \(S\) 仿射的情形.
由 3.5, 此时 \(G\) 有一个拟紧开子集 \(U\), 包含
\[
  \underline{G^0}=G.
\]
因此 \(G\) 拟紧, 因而在仿射概形 \(S\) 上拟紧 (EGA~I, 6.6.4(v)).
% Translation statistics: source_lines=264; translated_lines=248; inline_math=134; display_math=23; pending_TODO=0.
% Review: complete component translation; math fragments, reference keys, environments and note counts checked against source; formula prose translated; compilation pending parent integration.
